\documentclass[11pt,a4paper]{article}

\usepackage[utf8]{inputenc}
\usepackage[T1]{fontenc}
\usepackage{lmodern}
\usepackage[margin=1in]{geometry}
\usepackage{amsmath,amssymb}
\usepackage{enumitem}
\usepackage{booktabs}
\usepackage{fancyhdr}
\usepackage{xcolor}
\usepackage{mdframed}

% Gray solution boxes
\definecolor{solngray}{gray}{0.95}
\newmdenv[backgroundcolor=solngray,linewidth=0.5pt,linecolor=gray,innertopmargin=8pt,innerbottommargin=8pt]{solnbox}

\pagestyle{fancy}
\fancyhf{}
\fancyhead[L]{\textbf{Discrete Structures I}}
\fancyhead[R]{\textbf{Prelim Practice Exam C}}
\fancyfoot[C]{\thepage}
\renewcommand{\headrulewidth}{0.5pt}

\begin{document}

\begin{center}
{\LARGE\textbf{DISCRETE STRUCTURES I}}\\[5pt]
{\Large Prelim Practice Exam --- Version C}\\[5pt]
{\large\textit{With Complete Solutions}}\\[10pt]
\rule{0.8\textwidth}{0.5pt}\\[5pt]
\textbf{Time Allowed:} 1.5 hours \quad | \quad \textbf{Total Points:} 55
\end{center}

\vspace{10pt}

\noindent\textbf{Instructions:} Answer all questions. Show all work for full credit.

\vspace{15pt}

%==============================================================================
% PART I: MULTIPLE CHOICE
%==============================================================================
\section*{PART I: Multiple Choice (10 points)}
\textit{2 points each.}

\vspace{8pt}

\textbf{1.} The implication $p \to q$ is false:

\begin{enumerate}[label=(\Alph*),noitemsep]
    \item When $p$ is false and $q$ is true
    \item When $p$ is true and $q$ is false
    \item When both $p$ and $q$ are false
    \item When both $p$ and $q$ are true
\end{enumerate}

\begin{solnbox}
\textbf{Answer: (B)}

\textbf{Explanation:} This is THE most important thing to memorize about implications!

The implication $p \to q$ is FALSE \textbf{only} when $p$ is TRUE and $q$ is FALSE.

\textbf{Think of it as a promise:}
\begin{itemize}[noitemsep]
    \item ``If it rains, I'll bring an umbrella.''
    \item Rain + umbrella = promise kept (T $\to$ T = T)
    \item Rain + no umbrella = \textbf{promise broken} (T $\to$ F = F)
    \item No rain + umbrella = promise not tested, so not broken (F $\to$ T = T)
    \item No rain + no umbrella = promise not tested, so not broken (F $\to$ F = T)
\end{itemize}

\textbf{Key insight:} A false hypothesis makes the implication trivially true (``vacuously true'').
\end{solnbox}

\vspace{10pt}

\textbf{2.} De Morgan's Law states that $\neg(p \land q)$ is equivalent to:

\begin{enumerate}[label=(\Alph*),noitemsep]
    \item $\neg p \land \neg q$
    \item $\neg p \lor \neg q$
    \item $p \lor q$
    \item $\neg p \to \neg q$
\end{enumerate}

\begin{solnbox}
\textbf{Answer: (B)}

\textbf{De Morgan's Laws:}
$$\neg(p \land q) \equiv \neg p \lor \neg q$$
$$\neg(p \lor q) \equiv \neg p \land \neg q$$

\textbf{Memory trick:} ``Break the bar, change the sign'' or ``NOT distributes and flips AND/OR.''

When you negate a compound statement:
\begin{enumerate}[noitemsep]
    \item Negate each component
    \item Flip AND to OR (or OR to AND)
\end{enumerate}

\textbf{Intuition:} For the conjunction $p \land q$ to be false (i.e., $\neg(p \land q)$), we need at least one of them to be false, i.e., $\neg p$ OR $\neg q$.
\end{solnbox}

\vspace{10pt}

\textbf{3.} Given the statement ``All cats are mammals,'' which is its contrapositive?

\begin{enumerate}[label=(\Alph*),noitemsep]
    \item If something is not a cat, then it is not a mammal.
    \item If something is a mammal, then it is a cat.
    \item If something is not a mammal, then it is not a cat.
    \item Some cats are not mammals.
\end{enumerate}

\begin{solnbox}
\textbf{Answer: (C)}

\textbf{Step 1:} Translate the original to logic.

``All cats are mammals'' = $\forall x (Cat(x) \to Mammal(x))$

For a specific $x$: $C(x) \to M(x)$

\textbf{Step 2:} Form the contrapositive.

Contrapositive of $p \to q$ is $\neg q \to \neg p$.

So contrapositive of $C(x) \to M(x)$ is $\neg M(x) \to \neg C(x)$.

\textbf{Step 3:} Translate back to English.

``If something is NOT a mammal, then it is NOT a cat.''

This is option (C)!

\textbf{Note:}
\begin{itemize}[noitemsep]
    \item (A) is the INVERSE (negate both, but don't swap) --- not equivalent
    \item (B) is the CONVERSE (swap, but don't negate) --- not equivalent
    \item (C) is the CONTRAPOSITIVE --- equivalent!
\end{itemize}
\end{solnbox}

\vspace{10pt}

\textbf{4.} $\forall x \exists y\, P(x, y)$ differs from $\exists y \forall x\, P(x, y)$ in that:

\begin{enumerate}[label=(\Alph*),noitemsep]
    \item The first says every $x$ has its own $y$; the second says one $y$ works for all $x$
    \item They are logically equivalent
    \item The first is always false
    \item The first uses universal, the second uses existential quantification only
\end{enumerate}

\begin{solnbox}
\textbf{Answer: (A)}

\textbf{This is crucial! Quantifier order matters when quantifiers are different!}

\textbf{$\forall x \exists y\, P(x, y)$:} ``For every $x$, there exists some $y$ such that $P(x, y)$.''

The $y$ can depend on $x$. Each $x$ gets its own $y$.

Example: ``Everyone has a mother.'' Each person has their own mother.

\textbf{$\exists y \forall x\, P(x, y)$:} ``There exists a $y$ such that for all $x$, $P(x, y)$.''

ONE fixed $y$ must work for ALL $x$.

Example: ``There is someone who is everyone's friend.'' One universal friend.

\textbf{These are NOT equivalent!} The second is a stronger claim than the first.
\end{solnbox}

\vspace{10pt}

\textbf{5.} The argument ``If I study, I pass. I did not study. Therefore, I did not pass'' is:

\begin{enumerate}[label=(\Alph*),noitemsep]
    \item Valid by Modus Tollens
    \item Valid by Modus Ponens
    \item Invalid (Denying the Antecedent)
    \item Invalid (Affirming the Consequent)
\end{enumerate}

\begin{solnbox}
\textbf{Answer: (C)}

Let $s$ = ``I study'', $p$ = ``I pass''

The argument is: $s \to p$, $\neg s$, therefore $\neg p$.

This is the \textbf{Denying the Antecedent} fallacy!

\textbf{Why it's invalid:}

Just because studying leads to passing doesn't mean it's the ONLY way to pass!

Maybe you passed by:
\begin{itemize}[noitemsep]
    \item Lucky guessing
    \item Already knowing the material
    \item The test being curved
\end{itemize}

\textbf{Compare with valid forms:}
\begin{itemize}[noitemsep]
    \item Modus Ponens: $s \to p$, $s$ $\therefore p$ (valid)
    \item Modus Tollens: $s \to p$, $\neg p$ $\therefore \neg s$ (valid)
    \item Denying Antecedent: $s \to p$, $\neg s$ $\therefore \neg p$ (INVALID!)
    \item Affirming Consequent: $s \to p$, $p$ $\therefore s$ (INVALID!)
\end{itemize}
\end{solnbox}

\newpage

%==============================================================================
% PART II: SHORT ANSWER
%==============================================================================
\section*{PART II: Short Answer (20 points)}

\textbf{Question 1.} (5 points) Construct a truth table for $(p \oplus q) \leftrightarrow \neg(p \leftrightarrow q)$.

\begin{solnbox}
\textbf{This problem reveals an important relationship!}

$\oplus$ is XOR (exclusive or): true when exactly one operand is true.

$\leftrightarrow$ is biconditional (iff): true when both operands have the same value.

\textbf{Prediction:} XOR is the ``opposite'' of biconditional, so $p \oplus q \equiv \neg(p \leftrightarrow q)$.

\begin{center}
\begin{tabular}{cc|cc|c|c}
\toprule
$p$ & $q$ & $p \oplus q$ & $p \leftrightarrow q$ & $\neg(p \leftrightarrow q)$ & $(p \oplus q) \leftrightarrow \neg(p \leftrightarrow q)$ \\
\midrule
T & T & F & T & F & T \\
T & F & T & F & T & T \\
F & T & T & F & T & T \\
F & F & F & T & F & T \\
\bottomrule
\end{tabular}
\end{center}

\textbf{Result:} All T's in the final column! This is a \textbf{TAUTOLOGY}.

\textbf{What this means:} $p \oplus q \equiv \neg(p \leftrightarrow q)$

XOR and biconditional are negations of each other:
\begin{itemize}[noitemsep]
    \item $\leftrightarrow$: same truth values $\to$ true
    \item $\oplus$: different truth values $\to$ true
\end{itemize}
\end{solnbox}

\vspace{15pt}

\textbf{Question 2.} (5 points) Translate into propositional logic:

\textit{``A user can delete a file only if they have write permissions or they are the administrator.''}

Let:
\begin{itemize}[noitemsep]
    \item $d$: ``User can delete file''
    \item $w$: ``User has write permissions''
    \item $a$: ``User is administrator''
\end{itemize}

\begin{solnbox}
\textbf{Key phrase:} ``$A$ only if $B$'' means $A \to B$.

So: ``User can delete only if [condition]'' = $d \to [\text{condition}]$

The condition is: ``write permissions OR administrator'' = $w \lor a$

\textbf{Answer:}
$$\boxed{d \to (w \lor a)}$$

\textbf{Interpretation:} Having write permissions or being admin is NECESSARY for deletion. It's a requirement, not a guarantee.

\textbf{Note:} This doesn't say deleting is sufficient for having permissions. It says if you CAN delete, you MUST have one of those permissions.
\end{solnbox}

\vspace{15pt}

\textbf{Question 3.} (4 points) Determine truth values. Domain: all integers.

\begin{enumerate}[label=(\roman*)]
    \item $\exists x \forall y\, (x + y = y)$
    \item $\forall x \exists y\, (x \cdot y = 1)$
\end{enumerate}

\begin{solnbox}
\textbf{(i) $\exists x \forall y\, (x + y = y)$}

``There exists an integer $x$ such that for ALL integers $y$, $x + y = y$.''

We need ONE $x$ that works universally.

$x + y = y$ means $x = 0$.

Check: $0 + y = y$ for all $y$? Yes! Zero is the additive identity.

This is \textbf{TRUE}. The witness is $x = 0$.

\rule{\linewidth}{0.3pt}

\textbf{(ii) $\forall x \exists y\, (x \cdot y = 1)$}

``For every integer $x$, there exists an integer $y$ such that $x \cdot y = 1$.''

We need every integer to have a multiplicative inverse \textit{that is also an integer}.

Consider $x = 2$: We need $2y = 1$, so $y = 1/2$. But $1/2$ is NOT an integer!

Consider $x = 0$: We need $0 \cdot y = 1$, but $0 \cdot y = 0$ for all $y$. Impossible!

This is \textbf{FALSE}.

\textbf{Note:} If the domain were REAL numbers (except 0), this would be true.
\end{solnbox}

\vspace{15pt}

\textbf{Question 4.} (6 points) Let $T(x, y)$ mean ``$x$ is taller than $y$'' where the domain is all people. Express:

\begin{enumerate}[label=(\roman*)]
    \item Everyone is taller than someone.
    \item There is someone who is not taller than anyone.
    \item Someone is the tallest (taller than everyone else).
\end{enumerate}

\begin{solnbox}
\textbf{(i) Everyone is taller than someone.}

``For every person $x$, there exists a person $y$ such that $x$ is taller than $y$.''

$$\boxed{\forall x \exists y\, T(x, y)}$$

(Each person has someone shorter than them.)

\rule{\linewidth}{0.3pt}

\textbf{(ii) There is someone who is not taller than anyone.}

``There exists a person $x$ such that for all people $y$, $x$ is NOT taller than $y$.''

$$\boxed{\exists x \forall y\, \neg T(x, y)}$$

(This describes the shortest person --- they're not taller than anyone!)

\rule{\linewidth}{0.3pt}

\textbf{(iii) Someone is the tallest.}

``There exists a person $x$ such that for all OTHER people $y$, $x$ is taller than $y$.''

$$\boxed{\exists x \forall y\, (x \neq y \to T(x, y))}$$

\textbf{Why the $x \neq y$?} We don't want to require someone to be taller than themselves (which is impossible).

Alternative (simpler, if we allow $T(x,x)$ to be false naturally):
$\exists x \forall y\, (y \neq x \to T(x, y))$
\end{solnbox}

\newpage

%==============================================================================
% PART III: PROBLEM SOLVING
%==============================================================================
\section*{PART III: Problem Solving (25 points)}

\textbf{Question 5.} (6 points) Identify the rule of inference:

\begin{enumerate}[label=(\roman*)]
    \item I will go to the party and I will dance. Therefore, I will dance.
    \item If $x$ is even, then $x^2$ is even. $x = 4$ is even. Therefore, $16$ is even.
    \item If the alarm sounds, there is danger or a test is happening. The alarm sounds. Therefore, there is danger or a test is happening.
\end{enumerate}

\begin{solnbox}
\textbf{(i)} $p \land q$, therefore $q$

This is \textbf{Simplification}!

If a conjunction is true, each conjunct is individually true.

\rule{\linewidth}{0.3pt}

\textbf{(ii)} Let $E(x)$ = ``$x$ is even'', consider specific value $x = 4$.

Structure:
\begin{align*}
&\forall x\, (E(x) \to E(x^2)) \\
&E(4) \\
&\therefore E(16)
\end{align*}

This uses TWO rules:
\begin{enumerate}[noitemsep]
    \item \textbf{Universal Instantiation}: From $\forall x\, P(x)$ conclude $P(4)$: ``If 4 is even, then $4^2$ is even''
    \item \textbf{Modus Ponens}: From ``If 4 is even, then 16 is even'' and ``4 is even'', conclude ``16 is even''
\end{enumerate}

\rule{\linewidth}{0.3pt}

\textbf{(iii)} Let $a$ = alarm, $d$ = danger, $t$ = test

Structure: $a \to (d \lor t)$, $a$, therefore $d \lor t$

This is \textbf{Modus Ponens}!

$$\frac{p \to q \quad p}{\therefore q}$$

The antecedent ($a$) is true, so the consequent ($d \lor t$) follows.
\end{solnbox}

\vspace{15pt}

\textbf{Question 6.} (7 points) Prove: $\neg p \lor (p \land q) \equiv p \to q$ (without truth tables).

\begin{solnbox}
\textbf{Strategy:} Transform the left side step by step.

\begin{align*}
\neg p \lor (p \land q) &\equiv (\neg p \lor p) \land (\neg p \lor q) && \text{[Distributive Law]} \\
&\equiv T \land (\neg p \lor q) && \text{[Negation Law: } \neg p \lor p = T \text{]} \\
&\equiv \neg p \lor q && \text{[Identity Law: } T \land X = X \text{]} \\
&\equiv p \to q && \text{[Implication Law]}
\end{align*}

$$\boxed{\neg p \lor (p \land q) \equiv p \to q}$$

\textbf{Laws used:}
\begin{enumerate}[noitemsep]
    \item Distributive Law: $A \lor (B \land C) \equiv (A \lor B) \land (A \lor C)$
    \item Negation Law: $\neg p \lor p \equiv T$
    \item Identity Law: $T \land X \equiv X$
    \item Implication Law: $\neg p \lor q \equiv p \to q$
\end{enumerate}

\textbf{Alternative approach} (going right to left):

Start with $p \to q \equiv \neg p \lor q$ and show it equals $\neg p \lor (p \land q)$.

By absorption: $\neg p \lor (p \land q) \equiv \neg p \lor q$ when $\neg p$ ``absorbs'' the $p$ in the second term.
\end{solnbox}

\vspace{15pt}

\textbf{Question 7.} (4 points) Negate and simplify:

``Every student who studies hard passes the exam.''

Express using $S(x)$ = ``$x$ is a student'', $H(x)$ = ``$x$ studies hard'', $P(x)$ = ``$x$ passes''.

\begin{solnbox}
\textbf{Step 1:} Translate the original.

``Every student who studies hard passes'' = ``For all $x$, if $x$ is a student and studies hard, then $x$ passes.''

$$\forall x\, ((S(x) \land H(x)) \to P(x))$$

\textbf{Step 2:} Negate it.

\begin{align*}
&\neg \forall x\, ((S(x) \land H(x)) \to P(x)) \\
&\equiv \exists x\, \neg((S(x) \land H(x)) \to P(x)) && \text{[De Morgan for quantifiers]}
\end{align*}

\textbf{Step 3:} Simplify $\neg(A \to B)$.

Recall: $\neg(A \to B) \equiv A \land \neg B$

So: $\neg((S(x) \land H(x)) \to P(x)) \equiv (S(x) \land H(x)) \land \neg P(x)$

\textbf{Final answer:}
$$\boxed{\exists x\, (S(x) \land H(x) \land \neg P(x))}$$

\textbf{In English:} ``There exists a student who studies hard but does NOT pass the exam.''

\textbf{Intuition:} The negation of ``all such students pass'' is ``at least one such student doesn't pass.''
\end{solnbox}

\vspace{15pt}

\textbf{Question 8.} (8 points) Prove valid:

\textbf{Premises:}
\begin{enumerate}[noitemsep]
    \item $(p \land q) \to r$
    \item $p$
    \item $q$
\end{enumerate}

\textbf{Conclusion:} $r$

\begin{solnbox}
\textbf{Strategy:} Combine $p$ and $q$, then apply Modus Ponens.

\begin{center}
\begin{tabular}{c|l|l}
\toprule
\textbf{Step} & \textbf{Statement} & \textbf{Reason} \\
\midrule
1 & $p$ & Premise \\
2 & $q$ & Premise \\
3 & $p \land q$ & Conjunction (1, 2) \\
4 & $(p \land q) \to r$ & Premise \\
5 & $r$ & Modus Ponens (3, 4) \\
\bottomrule
\end{tabular}
\end{center}

$$\boxed{\therefore r}$$

\textbf{Explanation:}
\begin{itemize}[noitemsep]
    \item Steps 1-3: We're given $p$ and $q$ separately. The Conjunction rule lets us combine them into $p \land q$.
    \item Steps 4-5: Now we have both $p \land q$ and $(p \land q) \to r$. The hypothesis of the implication is satisfied, so by Modus Ponens, the conclusion $r$ follows.
\end{itemize}

\textbf{This is a very direct proof!} The key insight is recognizing you need to BUILD $p \land q$ from separate premises before you can use the conditional.
\end{solnbox}

\vspace{20pt}

\begin{center}
\rule{0.5\textwidth}{0.5pt}\\[5pt]
\textbf{END OF EXAM}\\[3pt]
\textit{Total: 55 points}
\end{center}

\end{document}
